\documentclass[11pt]{article}
\usepackage[margin=1in]{geometry}
\usepackage{booktabs}
\usepackage{amsmath}
\usepackage{graphicx}
\usepackage{xcolor}
\usepackage{url}
\usepackage[hidelinks]{hyperref}
\usepackage{caption}

\title{Are Agent Memory Systems Worth It for Coding Agents?\\
\large MemBench: A Benchmark-Harness Pilot with Pinned Open-Weight Models and Explicit Contamination Probes}
\author{MemBench Project\thanks{Benchmark code, dataset, configurations, and full agent trajectories at the MemBench repository. Inference: Qwen3-Coder-30B (AWQ 4-bit) served by vLLM 0.24.0 on a single RTX 5090.}}
\date{July 3, 2026}

\begin{document}
\maketitle

\begin{abstract}
Persistent memory systems for LLM coding agents are widely marketed but weakly evidenced: every prominent published comparison is vendor-produced, and the two head-to-head disputes on record (Zep vs.\ Mem0 on LoCoMo) show that configuration choices alone swing reported scores by more than 25 percentage points. We present \textbf{MemBench}, an executable benchmark harness --- with explicit contamination \emph{probes} rather than a contamination guarantee --- for measuring whether agent memory earns its cost on coding tasks, and report a pilot study comparing \textbf{seven memory arms}, presented by mechanism with the product that inspired each: no-memory, raw-history RAG, structured markdown memory (Cline-Memory-Bank/Cursor-rules style), FTS/BM25 observations (claude-mem-style), vector-extraction memory (Mem0 OSS), graph-edge memory (Graphiti/FalkorDB), and KG-query memory (graphify) --- under a single agent scaffold (Claude Code) driven end-to-end by a \emph{pinned open-weight model} (Qwen3-Coder-30B via vLLM), with extraction and embedding for every memory system routed to the same local endpoints. Because Qwen3-Coder-30B (July 2025) postdates SWE-bench-Live's April-2025 freshness window, we do not claim contamination control for this backbone; a memorization probe and date-split analysis are the mitigations we report instead. On our five-task pilot suite, \emph{every arm resolved every task} (35/35, independently re-verified). The pilot's chief findings are methodological: (i) the agent scaffold dominates the model --- the same model scored 0/6 single-shot but 100\% agentic; (ii) un-isolated harnesses silently corrupt results --- before environment isolation, headless agents inherited user-level tools and auto-memory, token costs spread $2.4\times$ across arms, and the platform systems appeared most expensive; after isolation the spread collapsed to $1.36\times$ and the cost ranking \emph{inverted}, with Mem0 and graphify cheapest; and (iii) memory's first measurable effect is on agent \emph{behavior and cost}, not correctness. We release the harness, tasks, per-turn trajectories, and all failure modes as a foundation for the statistically powered evaluation the field currently lacks.
\end{abstract}

\section{Introduction}

Memory layers for coding agents promise cross-session learning: remember the project's conventions, past bugs, and setup quirks, and the agent should resolve issues faster and cheaper. The market took the promise seriously --- Mem0 raised \$24M, Zep commercialized Graphiti, and the most popular Claude Code memory plugin (claude-mem) exceeds 85k GitHub stars. The evidence did not keep pace:

\begin{itemize}
  \item \textbf{Vendor numbers do not survive audit.} Zep reported 84\% on LoCoMo; a corrected third-party evaluation measured $58.4\%$. Mem0's headline +26\% is measured on a benchmark whose answer key contains 6.4\% errors and whose LLM judge accepts 62.8\% of intentionally wrong answers~\cite{locomo-audit}.
  \item \textbf{Simple baselines embarrass platforms.} A plain filesystem-plus-grep agent beat Mem0's graph variant on Letta's own benchmark (74.0 vs.\ 68.5)~\cite{letta-fs}; on MemoryAgentBench, Mem0 scored 3.4\% on test-time learning against 87.6\% for a long-context baseline~\cite{memoryagentbench}.
  \item \textbf{The standard coding benchmark is compromised.} OpenAI's 2026 audit of SWE-bench Verified found 59.4\% flawed tests among consistent failures and verbatim gold-patch memorization across frontier models~\cite{swebench-audit}, making ``memory lift'' on it indistinguishable from contamination recall.
\end{itemize}

No coding-agent memory benchmark with executable oracles, contamination controls, pinned open-weight models, and statistical power exists. MemBench is our attempt to build one. This paper reports the design and a \emph{pilot}: a deliberately small, fully instrumented dry run whose purpose is to validate the harness end-to-end and surface the failure modes that would corrupt a larger study. We are explicit about scale: five tasks, one seed, one model --- the pilot answers ``does the methodology work?'', not yet ``which memory system wins?''.

\section{Related Work}

\textbf{Memory systems.} Mem0 (vector store with LLM fact extraction; the 2026 OSS release replaced the published two-phase pipeline with single-pass ADD-only extraction), Graphiti/Zep (bi-temporal knowledge graph, hybrid semantic+BM25+graph retrieval), Letta/MemGPT (self-edited memory blocks), Cognee (ECL knowledge-graph pipeline), claude-mem (Claude Code plugin: hook-captured observations in SQLite FTS5), graphify (repository knowledge graphs via tree-sitter + LLM semantic extraction), and Anthropic's native memory tool and auto-memory. Of these, only ReasoningBank-style distilled episodic memory has a measured lift on a repository-level coding benchmark (+4.6pp on SWE-bench Verified, with the contamination caveat above)~\cite{reasoningbank}.

\textbf{Benchmarks.} LoCoMo and LongMemEval (conversational recall; LoCoMo demonstrably unreliable), MemoryAgentBench (four memory competencies; conversational), EvoMemBench (memory value shrinks as context windows grow; its context-size curve bounds how much any memory arm can be worth under a large window), STALE (best frontier model detects invalidated memories only 55.2\% of the time), and SWE Context Bench (319 interdependent GitHub tasks; oracle summary reuse +8pp). The closest prior work is \textbf{SWE-Bench-CL}~\cite{swebenchcl}: chronological per-repo task sequences on SWE-bench Verified evaluated with a single fixed FAISS semantic-memory agent. MemBench differs by (i) a post-2024 SWE-bench-Live base rather than the contamination-exposed Verified set, (ii) seven standardized memory arms rather than one fixed memory implementation, (iii) an explicit carry-forward control, and (iv) per-arm cost accounting. Concurrent coding-memory work we position against includes Memory Transfer Learning across coding agents~\cite{memtransfer}, AgentCL~\cite{agentcl}, structurally aligned subtask-level memory~\cite{subtaskmem}, SlopCodeBench~\cite{slopcode}, BEAM~\cite{beam}, and the RL memory-policy class typified by Memory-R1~\cite{memoryr1}; the field-level framing follows the \emph{Memory in the Age of AI Agents} survey~\cite{memsurvey}. Two memory-mechanism classes our taxonomy adds are personalized-PageRank retrieval (HippoRAG / HippoRAG 2~\cite{hipporag,hipporag2}, which we flag as a planned neutral academic arm) and RL-learned memory policies (Memory-R1). For contamination-safe task sourcing we build on SWE-bench-Live (post-2024 issues, monthly refresh) and the SWE-bench harness pattern of Dockerized per-instance environments with fail-to-pass/pass-to-pass oracles~\cite{swebench-live}.

\textbf{Measurement hazards.} Run-to-run pass@1 on SWE-bench-style evals varies 2.2--6.0pp across identical reruns --- the same order as typical memory effect sizes --- so single-run comparisons are uninterpretable~\cite{errorbars}. Our companion design document additionally identifies a \emph{carry-forward hazard} specific to memory evaluation: in 70\% of consecutive Django task pairs, the earlier task's fix is already present in the later task's base checkout, inflating naive ``memory lift''.

\section{MemBench Design}

\subsection{Task model}
Each instance models an agent returning to a project: an issue, a workspace snapshot (template repository), an executable oracle (pytest fail-to-pass and pass-to-pass sets, with test paths protected from agent modification), and a time-split \emph{memory corpus} --- structured project memory, event history, and documents, all strictly predating the task. A pre-check runs the fail-to-pass tests before the agent starts; an instance only counts as resolved if the bug demonstrably existed, the fail-to-pass tests pass afterward, and no pass-to-pass test regresses.

\subsection{Pilot task suite}
Five small Python tasks instantiate the design's memory-need taxonomy:
\begin{center}
\begin{tabular}{lll}
\toprule
Instance & Category & Description \\
\midrule
demo\_\_001 & M2 episodic & Decimal division leaks backend exception \\
demo\_\_002 & M2 repeated pattern & Same leak class, different module; past fix in corpus \\
demo\_\_003 & M3 convention & Legacy config aliases; convention only in memory \\
demo\_\_004 & S stale trap & Corpus contains an invalidated module-location note \\
demo\_\_005 & M0 + distractors & Trivial off-by-one; corpus is all distractors \\
\bottomrule
\end{tabular}
\end{center}

\subsection{Memory arms}
\label{subsec:arms}
All arms share one scaffold --- Claude Code in headless mode --- and differ only in what the memory layer injects into the task prompt. Retrieved items are capped at 2{,}000 tokens.

\begin{center}
\small
\begin{tabular}{llll}
\toprule
Arm & Mechanism (product inspiration) & Storage / retrieval & Extraction LLM \\
\midrule
none & no memory & --- & --- \\
raw\_rag & raw-history RAG & corpus JSONL, lexical scoring & none \\
structured & structured markdown memory (Cline-Memory-Bank / Cursor-rules) & curated facts, lexical & none \\
claude\_mem & FTS/BM25 observations (claude-mem-style) & SQLite FTS5, BM25 & none (corpus pre-structured) \\
mem0 & vector-extraction memory (Mem0 OSS 2.0.11) & Qdrant (local), vector & qwen3-coder-30b (vLLM provider) \\
graphiti & graph-edge memory (Graphiti 0.29.x) & FalkorDB temporal graph, hybrid & qwen3-coder-30b + qwen3-embedding \\
graphify & KG-query memory (graphify) & knowledge graph, BFS query & qwen3-coder-30b (OpenAI-compat) \\
\bottomrule
\end{tabular}
\end{center}

The claude-mem arm replicates the plugin's storage and search layer (FTS5, BM25 ranking) rather than running the shipped plugin: the installed release cannot re-point its extraction step at a local model, and running it against a commercial API would break the pinned-model control. Letta and Cognee are excluded from the pilot: Letta's Claude Code bridge requires its cloud service, and Cognee's documentation recommends larger extraction models than our pinned backbone. Crucially, the product arms (Mem0, Graphiti, claude-mem) run in a \emph{standardized corpus-ingestion mode}: the driver hands each system the same time-split corpus and calls its retrieval interface, but the products' native incremental write and consolidation pipelines --- the mechanisms that would run during live deployment --- do not execute. These arms therefore measure \emph{retrieval mechanisms over a driver-authored store}, not the products as deployed. One consequence is a cost artifact: arms that re-ingest the growing corpus per task incur an $O(k^2)$ re-ingestion overhead across a $k$-task sequence that a live incremental store would not pay; reported cost numbers for the caching arms include this artifact.

\subsection{Pinned open-weight stack}
Reproducibility requires that no arm silently depends on a commercial API. The agent, and the extraction/embedding calls of every memory system, run against two local endpoints: Qwen3-Coder-30B-A3B-Instruct (AWQ 4-bit) and Qwen3-Embedding-0.6B, both served by vLLM 0.24.0 on one RTX 5090 (32\,GB). vLLM's native Anthropic-compatible \texttt{/v1/messages} endpoint lets Claude Code drive the open-weight model directly; tool calling requires \texttt{--enable-auto-tool-choice --tool-call-parser qwen3\_coder}.

\subsection{Metrics}
Per instance and arm: resolution (executable oracle only --- no LLM judges), wall time, API turns, per-tool call counts (from full stream-JSON trajectories), total tokens, retrieved memory items, and memory-store construction time. Cost for the pinned stack is electricity; token counts are the portable cost signal. Beyond the pilot's saturated single-cell counts, the harness now (concurrently implemented) reports a $k{\geq}2$ headline resolution split (memory value is only testable once a repo sequence has history to carry), paired significance via McNemar on discordant pairs and paired bootstrap CIs, a \texttt{citation\_rate}/memory-utilization signal parsed from trajectories (did the agent actually read injected memory?), an injected-token covariate per arm (to separate ``more context'' from ``better memory''), a date-split by issue creation month, and stratification by carry-forward status and offline memory-need label. A memorization probe (patch requested with no repo access) is reported per instance. We state status honestly: resolution, cost, turn, tool-count, and construction-time metrics are implemented and reported in this pilot; the $k{\geq}2$ split, McNemar/bootstrap CIs, citation-rate parsing, injected-token covariate, date-split, carry-forward/memory-need stratification, and memorization probe are implemented for the funded next step but not exercised on the five-task pilot suite, which lacks the repo history and cell counts to populate them.

\section{Harness Findings: What Silently Corrupts Memory Benchmarks}
\label{sec:harness}

The pilot's most transferable results are five failure modes we hit, diagnosed from trajectories, and fixed. Any of them would have invalidated a larger study.

\begin{enumerate}
  \item \textbf{Scaffold dominates model.} Single-shot patch generation with Qwen3-Coder-30B resolved 0/6 attempts (it cannot see the repository); the same model inside the agentic scaffold resolved every pilot task. Memory comparisons made on a crippled scaffold measure nothing about memory.
  \item \textbf{Headless agents never compact.} Claude Code assumes a 200k context for unknown models and disables auto-compaction in print mode; agents spiraled to 70+ turns and died at the 65k serving limit ($\sim$7 minutes and 3.4M input tokens per failure). Fixes: declare the real window (\texttt{CLAUDE\_CODE\_MAX\_CONTEXT\_TOKENS}), cap output tokens, and impose an explicit \texttt{--max-turns} budget --- itself a benchmark parameter.
  \item \textbf{Environment contamination.} Un-isolated headless runs inherit the host user's Claude Code configuration: task-management tools, installed skills, third-party memory plugins (including claude-mem's capture hooks), and auto-memory, which wrote \texttt{MEMORY.md} files that appeared \emph{inside model patches}. Trajectories showed agents spending 60--80\% of capped turns on inherited task tools after the fix had landed. Benchmark runs must use an isolated configuration directory with auto-memory disabled.
  \item \textbf{Bounded failure is not harness error.} Hitting a turn budget must be recorded as a scoreable unresolved-or-resolved attempt (the oracle decides), not discarded: in \emph{every} capped pilot run the fix had already landed by turn $\sim$6--12, so discarding capped runs would have systematically undercounted resolutions.
  \item \textbf{Memory-platform pinning is fragile in practice.} Pinning all systems to one local stack surfaced: an embedding-dimension parameter rejected by vLLM (Mem0), a search-API signature change between the documented and shipped versions (Mem0 2.x), a global telemetry store that deadlocks concurrent corpora (Mem0), and a deprecated default graph backend that creates no indexes at all (Graphiti/Kuzu; we substituted the supported FalkorDB). None of these appear in any vendor benchmark report.
\end{enumerate}

\section{Pilot Results}
\label{sec:results}

\begin{table}[h]
\centering
\caption{Isolated pilot run: 7 memory arms $\times$ 5 tasks, Qwen3-Coder-30B driving Claude Code (30-turn budget). \emph{Re-verified} = recorded patch re-applied to a pristine checkout re-passes the oracle. \emph{Capped} = runs hitting the turn budget (all still resolved). Averages over the arm's 5 runs.}
\label{tab:main}
\small
\begin{tabular}{lcccrrrr}
\toprule
Arm & Resolved & Re-verified & Capped & Wall (s) & Turns & Tokens & Mem items \\
\midrule
graphify        & 5/5 & 5/5 & 2 & 36 & 20.4 & 460{,}143 & 1.0 \\
mem0            & 5/5 & 5/5 & 2 & 36 & 20.8 & 480{,}847 & 8.4 \\
raw\_rag        & 5/5 & 5/5 & 3 & 48 & 21.8 & 504{,}060 & 1.8 \\
none            & 5/5 & 5/5 & 3 & 44 & 23.0 & 522{,}292 & 0.0 \\
structured      & 5/5 & 5/5 & 3 & 38 & 23.8 & 540{,}325 & 1.4 \\
graphiti        & 5/5 & 5/5 & 3 & 39 & 24.0 & 547{,}330 & 5.0 \\
claude\_mem     & 5/5 & 5/5 & 4 & 51 & 27.0 & 624{,}544 & 3.8 \\
\bottomrule
\end{tabular}
\end{table}

\subsection{Resolution and verification}
All seven arms resolved all five tasks (Table~\ref{tab:main}). Every resolution was verified outside the harness: the recorded source patch was applied to a pristine template checkout and the oracle re-run (35/35 re-pass). Oracle prechecks confirmed each bug failed before the agent ran, and trajectory analysis confirmed the fix landed well inside the turn budget in every capped run.

\subsection{Isolation flipped the cost ranking}
The identical matrix run \emph{without} environment isolation produced a different picture: token costs spread $2.4\times$ (561k--1.36M per task), the platform arms (Mem0, Graphiti) appeared most expensive and hit the turn cap most often, and trajectories showed agents spending most post-fix turns on tools inherited from the host user's environment. After isolation, the spread collapsed to $1.36\times$ (460k--625k) and the ranking inverted --- Mem0 and graphify became the \emph{cheapest} arms. The pre-isolation ``platform bloat'' was largely contamination, not the platforms. We consider this the pilot's most important quantitative result: an un-isolated harness did not merely add noise, it \emph{systematically reordered} the systems under test.

\subsection{What cost differences survive isolation}
The surviving spread is modest and, at $n{=}1$ per cell, indistinguishable from run-to-run variance (arms swung 0/5--5/5 across reruns during harness debugging when infrastructure faults intervened). Two qualitative observations still hold across all runs: injected memory count is uncorrelated with success on tasks this easy (Mem0's 8.4 items and none's 0 both score 5/5), and the distractor-corpus task (M0) shows memory arms paying for irrelevant context the no-memory arm never sees --- the overhead channel exists even when it does not change outcomes.

\subsection{Turn-budget behavior}
2--4 runs per arm hit the 30-turn budget, yet every capped run had already produced the fix (typically by turn 6--12); remaining turns went to redundant verification and task-management churn. The budget therefore acted as a pure cost control on this suite --- it changed no outcomes --- but it is a benchmark parameter that harder tasks will interact with, and must be reported.

\section{Discussion: Are Memory Systems Worth It?}

The pilot cannot answer the headline question --- its tasks are too easy (the no-memory arm scores 100\%) and $n{=}1$ per cell sits inside known run-to-run variance. What it does establish:

\begin{itemize}
  \item \textbf{Harness hygiene is worth more percentage points than memory system choice.} Environment isolation moved per-arm token costs by up to $2.8\times$ and reordered the systems under test; no published memory comparison we are aware of reports its isolation regime. Until a benchmark controls scaffold, budget, and environment, its ranking of memory systems is not measuring memory systems.
  \item \textbf{On tasks the agent can solve alone, memory is pure overhead --- but small.} With correctness saturated, the isolated cost spread across arms was 1.36$\times$; memory construction (milliseconds for FTS, seconds for graph extraction, longest for Mem0's per-item LLM pipeline) remains the larger practical difference at corpus scale.
  \item \textbf{Memory's first measurable effect is behavioral.} Injected context changed how long agents ran and what they did after fixing the bug, before it changed whether they fixed it. Cost and behavior metrics must be first-class in any memory evaluation.
  \item \textbf{A fair platform comparison on pinned open weights is feasible} --- all four integrated systems ran extraction and embedding on local endpoints --- but only after fixing failure modes that vendor-run evaluations have no incentive to report (Section~\ref{sec:harness}).
\end{itemize}

\section{Threats to Validity}
\label{sec:threats}

\paragraph{Contamination (construct/internal).} Qwen3-Coder-30B was released in July 2025 and therefore postdates SWE-bench-Live's April-2025 freshness window; the instance issues we draw from (roughly 2024-10 through 2025-03) plausibly lie inside its training data. We consequently make \emph{no} ``contamination-controlled'' claim for this backbone. The mitigations we report instead are a memorization probe (request the patch with no repository access, per the arXiv:2506.12286 method) and a date-split analysis; a truly contamination-safe run requires a model whose cutoff precedes the task window, or post-cutoff instances relative to the chosen model.

\paragraph{Repository scope (external).} The eval covers 6 of 39 eligible repositories, all Python, selected for fast Docker evaluation. Aggregate numbers can therefore hide per-repo structure, so per-repo reporting is mandatory. Sequence-depth claims are especially thin: results at depth $k{\geq}8$ rest on at most three repositories, and $k{\geq}12$ on a single haystack repository. Depth-conditioned conclusions should be read as suggestive, not established.

\paragraph{Harness tuning as an upper-bound regime (external).} The 64k-window harness is tuned for memory to look good: a compact context window, a 400-line read cap, and tool-output trimming, tuned on the briefcase and haystack repositories, which are 41\% of the eval set. Under a large context window memory has less to add (EvoMemBench's context-size curve, cited above), so our results should be read as an \emph{upper-bound regime} for memory value rather than a neutral operating point.

\paragraph{Sequence construction (construct).} Task sequences are chronological, not interdependence-mined; we do not claim each later task genuinely depends on an earlier one. To keep ``memory lift'' from being inflated by answers already sitting in the repository, carry-forward status and memory-need labels are computed offline and used to stratify the analysis, rather than assumed.

\paragraph{Products are not measured as deployed (construct).} As noted in Section~\ref{subsec:arms}, the product arms (Mem0, Graphiti, claude-mem) run in standardized corpus-ingestion mode; their native incremental write and consolidation pipelines do not execute. These arms measure retrieval mechanisms over a driver-authored store, and the cost figures for caching arms include an $O(k^2)$ re-ingestion artifact that a live incremental store would not incur. A verdict of ``product X is/is not worth it'' from these arms is a statement about its retrieval layer, not its shipped write/consolidation behavior.

\section{Limitations and Next Steps}
Five synthetic tasks, one seed, one ordering, one model, no memory-write evaluation, and a corpus small enough that construction-time differences (milliseconds to minutes) do not yet bind. The funded next step is the designed pilot: 100--150 SWE-bench-Live-derived tasks in dependency-mined sequences with carry-forward annotation, 4--5 seeds $\times$ 3 orderings, paired McNemar and clustered-bootstrap analysis, stale/distractor/poison splits, and the decision rule pre-registered in the design document: if simple structured memory is within noise of the platforms, the recommendation is files and grep.

\section*{Reproducibility}
All configurations, task templates, memory corpora, adapter implementations, run outputs, and full per-turn trajectories (stream-JSON, including every tool call) accompany the repository. The serving stack is two pinned open-weight checkpoints on one consumer GPU; total pilot cost was local electricity.

\begin{thebibliography}{99}
\bibitem{swebench-audit} OpenAI. \emph{Why we no longer evaluate SWE-bench Verified}. Feb 23, 2026.
\bibitem{locomo-audit} Mem0 vs.\ Zep LoCoMo dispute: corrected evaluation, GitHub issue thread, May 2025; Zep, \emph{Lies, Damn Lies, \& Statistics}, 2025.
\bibitem{letta-fs} Letta. \emph{Benchmarking AI agent memory: is a filesystem all you need?} 2025.
\bibitem{memoryagentbench} Hu et al. \emph{MemoryAgentBench}. ICLR 2026. arXiv:2507.05257.
\bibitem{reasoningbank} \emph{ReasoningBank}. arXiv:2509.25140.
\bibitem{swebench-live} \emph{SWE-bench-Live}. arXiv:2505.23419.
\bibitem{errorbars} Miller et al. \emph{Adding Error Bars to Evals}. arXiv:2411.00640.
\bibitem{stale} \emph{STALE: Testing memory invalidation}. arXiv:2605.06527.
\bibitem{swectx} \emph{SWE Context Bench}. arXiv:2602.08316.
\bibitem{swebenchcl} \emph{SWE-Bench-CL: Continual Learning for Coding Agents}. arXiv:2507.00014.
\bibitem{memtransfer} \emph{Memory Transfer Learning across Coding Agents}. arXiv:2604.14004.
\bibitem{agentcl} \emph{AgentCL: Continual Learning for LLM Agents}. arXiv:2606.02461.
\bibitem{subtaskmem} \emph{Structurally Aligned Subtask-Level Memory}. arXiv:2602.21611.
\bibitem{slopcode} \emph{SlopCodeBench}. arXiv:2603.24755.
\bibitem{beam} \emph{BEAM: Benchmarking memory for agents}.
\bibitem{hipporag} \emph{HippoRAG: Neurobiologically Inspired Long-Term Memory for LLMs}. arXiv:2405.14831.
\bibitem{hipporag2} \emph{HippoRAG 2}. arXiv:2502.14802.
\bibitem{memoryr1} \emph{Memory-R1: Learning Memory Policies with RL}. arXiv:2508.19828.
\bibitem{memsurvey} \emph{Memory in the Age of AI Agents}. arXiv:2512.13564.
\end{thebibliography}

\end{document}
